%%%%%%%%%%%%%%%%%%%%%%%%%%%%%%%%%%%%%%%%%%%%%%%%%%%%%%%%%%%%%%%%%%%%%%%%
%%%%%%%%%%%%%%%%%%%%%% Simple LaTeX CV Template %%%%%%%%%%%%%%%%%%%%%%%%
%%%%%%%%%%%%%%%%%%%%%%%%%%%%%%%%%%%%%%%%%%%%%%%%%%%%%%%%%%%%%%%%%%%%%%%%

%%%%%%%%%%%%%%%%%%%%%%%%%%%%%%%%%%%%%%%%%%%%%%%%%%%%%%%%%%%%%%%%%%%%%%%%
%% NOTE: If you find that it says                                     %%
%%                                                                    %%
%%                           1 of ??                                  %%
%%                                                                    %%
%% at the bottom of your first page, this means that the AUX file     %%
%% was not available when you ran LaTeX on this source. Simply RERUN  %%
%% LaTeX to get the ``??'' replaced with the number of the last page  %%
%% of the document. The AUX file will be generated on the first run   %%
%% of LaTeX and used on the second run to fill in all of the          %%
%% references.                                                        %%
%%%%%%%%%%%%%%%%%%%%%%%%%%%%%%%%%%%%%%%%%%%%%%%%%%%%%%%%%%%%%%%%%%%%%%%%

%%%%%%%%%%%%%%%%%%%%%%%%%%%% Document Setup %%%%%%%%%%%%%%%%%%%%%%%%%%%%

% Don't like 10pt? Try 11pt or 12pt
\documentclass[10pt]{article}

% The automated optical recognition software used to digitize resume
% information works best with fonts that do not have serifs. This
% command uses a sans serif font throughout. Uncomment both lines (or at
% least the second) to restore a Roman font (i.e., a font with serifs).
%\usepackage{times}
%\renewcommand{\familydefault}{\sfdefault}

% This is a helpful package that puts math inside length specifications
\usepackage{calc}
\usepackage{comment}
%\usepackage{kotex}
\usepackage[normalem]{ulem}

% Simpler bibsection for CV sections
% (thanks to natbib for inspiration)
\makeatletter
\newlength{\bibhang}
\setlength{\bibhang}{1em} %1em}
\newlength{\bibsep}
 {\@listi \global\bibsep\itemsep \global\advance\bibsep by\parsep}
\newenvironment{bibsection}%
        {\begin{enumerate}{}{%
%        {\begin{list}{}{%
       \setlength{\leftmargin}{\bibhang}%
       \setlength{\itemindent}{-\leftmargin}%
       \setlength{\itemsep}{\bibsep}%
       \setlength{\parsep}{\z@}%
        \setlength{\partopsep}{0pt}%
        \setlength{\topsep}{0pt}}}
        {\end{enumerate}\vspace{-.6\baselineskip}}
%        {\end{list}\vspace{-.6\baselineskip}}
\makeatother

% Layout: Puts the section titles on left side of page
\reversemarginpar

%
%         PAPER SIZE, PAGE NUMBER, AND DOCUMENT LAYOUT NOTES:
%
% The next \usepackage line changes the layout for CV style section
% headings as marginal notes. It also sets up the paper size as either
% letter or A4. By default, letter was used. If A4 paper is desired,
% comment out the letterpaper lines and uncomment the a4paper lines.
%
% As you can see, the margin widths and section title widths can be
% easily adjusted.
%
% ALSO: Notice that the includefoot option can be commented OUT in order
% to put the PAGE NUMBER *IN* the bottom margin. This will make the
% effective text area larger.
%
% IF YOU WISH TO REMOVE THE ``of LASTPAGE'' next to each page number,
% see the note about the +LP and -LP lines below. Comment out the +LP
% and uncomment the -LP.
%
% IF YOU WISH TO REMOVE PAGE NUMBERS, be sure that the includefoot line
% is uncommented and ALSO uncomment the \pagestyle{empty} a few lines
% below.
%

%% Use these lines for letter-sized paper
\usepackage[paper=letterpaper,
            %includefoot, % Uncomment to put page number above margin
            marginparwidth=1.2in,     % Length of section titles
            marginparsep=.05in,       % Space between titles and text
            margin=.95in,             % 0.95 inch margins
            includemp]{geometry}

%% Use these lines for A4-sized paper
%\usepackage[paper=a4paper,
%            %includefoot, % Uncomment to put page number above margin
%            marginparwidth=30.5mm,    % Length of section titles
%            marginparsep=1.5mm,       % Space between titles and text
%            margin=25mm,              % 25mm margins
%            includemp]{geometry}

%% More layout: Get rid of indenting throughout entire document
\setlength{\parindent}{0in}

\usepackage[shortlabels]{enumitem}

%% Reference the last page in the page number
%
% NOTE: comment the +LP line and uncomment the -LP line to have page
%       numbers without the ``of ##'' last page reference)
%
% NOTE: uncomment the \pagestyle{empty} line to get rid of all page
%       numbers (make sure includefoot is commented out above)
%
\usepackage{fancyhdr,lastpage}
\pagestyle{fancy}
%\pagestyle{empty}      % Uncomment this to get rid of page numbers
\fancyhf{}\renewcommand{\headrulewidth}{0pt}
\fancyfootoffset{\marginparsep+\marginparwidth}
\newlength{\footpageshift}
\setlength{\footpageshift}
          {0.5\textwidth+0.5\marginparsep+0.5\marginparwidth-2in}
\lfoot{\hspace{\footpageshift}%
       \parbox{4in}{\, \hfill %
                    \arabic{page} of \protect\pageref*{LastPage} % +LP
%                    \arabic{page}                               % -LP
                    \hfill \,}}

% Finally, give us PDF bookmarks
\usepackage{color,hyperref}
\definecolor{darkblue}{rgb}{0.0,0.0,0.3}
\hypersetup{colorlinks,breaklinks,
            linkcolor=darkblue,urlcolor=darkblue,
            anchorcolor=darkblue,citecolor=darkblue}

%%%%%%%%%%%%%%%%%%%%%%%% End Document Setup %%%%%%%%%%%%%%%%%%%%%%%%%%%%


%%%%%%%%%%%%%%%%%%%%%%%%%%% Helper Commands %%%%%%%%%%%%%%%%%%%%%%%%%%%%

% The title (name) with a horizontal rule under it
% (optional argument typesets an object right-justified across from name
%  as well)
%
% Usage: \makeheading{name}
%        OR
%        \makeheading[right_object]{name}
%
% Place at top of document. It should be the first thing.
% If ``right_object'' is provided in the square-braced optional
% argument, it will be right justified on the same line as ``name'' at
% the top of the CV. For example:
%
%       \makeheading[\emph{Curriculum vitae}]{Your Name}
%
% will put an emphasized ``Curriculum vitae'' at the top of the document
% as a title. Likewise, a picture could be included:
%
%   \makeheading[\includegraphics[height=1.5in]{my_picutre}]{Your Name}
%
% the picture will be flush right across from the name.
\newcommand{\makeheading}[2][]%
        {\hspace*{-\marginparsep minus \marginparwidth}%
         \begin{minipage}[t]{\textwidth+\marginparwidth+\marginparsep}%
             {\large \bfseries #2 \hfill #1}\\[-0.15\baselineskip]%
                 \rule{\columnwidth}{1pt}%
         \end{minipage}}

% The section headings
%
% Usage: \section{section name}
\renewcommand{\section}[1]{\pagebreak[3]%
    \hyphenpenalty=10000%
    \vspace{1.3\baselineskip}%
    \phantomsection\addcontentsline{toc}{section}{#1}%
    \noindent\llap{\scshape\smash{\parbox[t]{\marginparwidth}{\raggedright #1}}}%
    \vspace{-\baselineskip}\par}

% An itemize-style list with lots of space between items
\newenvironment{outerlist}[1][\enskip\textbullet]%
        {\begin{itemize}[#1,leftmargin=*]}{\end{itemize}%
         \vspace{-.6\baselineskip}}

% An environment IDENTICAL to outerlist that has better pre-list spacing
% when used as the first thing in a \section
\newenvironment{lonelist}[1][\enskip\textbullet]%
        {\begin{list}{#1}{%
        \setlength{\partopsep}{0pt}%
        \setlength{\topsep}{0pt}}}
        {\end{list}\vspace{-.6\baselineskip}}

% An itemize-style list with little space between items
\newenvironment{innerlist}[1][\enskip\textbullet]%
        {\begin{itemize}[#1,leftmargin=*,parsep=0pt,itemsep=0pt,topsep=0pt,partopsep=0pt]}
        {\end{itemize}}

% An environment IDENTICAL to innerlist that has better pre-list spacing
% when used as the first thing in a \section
\newenvironment{loneinnerlist}[1][\enskip\textbullet]%
        {\begin{itemize}[#1,leftmargin=*,parsep=0pt,itemsep=0pt,topsep=0pt,partopsep=0pt]}
        {\end{itemize}\vspace{-.6\baselineskip}}

% To add some paragraph space between lines.
% This also tells LaTeX to preferably break a page on one of these gaps
% if there is a needed pagebreak nearby.
\newcommand{\blankline}{\quad\pagebreak[3]}
\newcommand{\halfblankline}{\quad\vspace{-0.5\baselineskip}\pagebreak[3]}

% Uses hyperref to link DOI
\newcommand\doilink[1]{\href{http://dx.doi.org/#1}{#1}}
\newcommand\doi[1]{doi:\doilink{#1}}

% For \url{SOME_URL}, links SOME_URL to the url SOME_URL
\providecommand*\url[1]{\href{#1}{#1}}
% Same as above, but pretty-prints SOME_URL in teletype fixed-width font
\renewcommand*\url[1]{\href{#1}{\texttt{#1}}}

% For \email{ADDRESS}, links ADDRESS to the url mailto:ADDRESS
\providecommand*\email[1]{\href{mailto:#1}{#1}}
% Same as above, but pretty-prints ADDRESS in teletype fixed-width font
%\renewcommand*\email[1]{\href{mailto:#1}{\texttt{#1}}}

%\providecommand\BibTeX{{\rm B\kern-.05em{\sc i\kern-.025em b}\kern-.08em
%    T\kern-.1667em\lower.7ex\hbox{E}\kern-.125emX}}
%\providecommand\BibTeX{{\rm B\kern-.05em{\sc i\kern-.025em b}\kern-.08em
%    \TeX}}
\providecommand\BibTeX{{B\kern-.05em{\sc i\kern-.025em b}\kern-.08em
    \TeX}}
\providecommand\Matlab{\textsc{Matlab}}

%%%%%%%%%%%%%%%%%%%%%%%% End Helper Commands %%%%%%%%%%%%%%%%%%%%%%%%%%%

%%%%%%%%%%%%%%%%%%%%%%%%% Begin CV Document %%%%%%%%%%%%%%%%%%%%%%%%%%%%

\begin{document}
\makeheading{Hyuk Jun Kweon}

\section{Contact Information}

% NOTE: Mind where the & separators and \\ breaks are in the following
%       table.
%
% ALSO: \rcollength is the width of the right column of the table
%       (adjust it to your liking; default is 1.85in).
%
\newlength{\rcollength}\setlength{\rcollength}{1.4in}%
%
\begin{tabular}[t]{@{}p{\textwidth-\rcollength-2\tabcolsep}p{\rcollength}@{}}
%\href{http://www.cse.osu.edu/}%
%     {Department of Computer Science and Engineering} & \\
%     \href{http://www.osu.edu/}{The Ohio State University}


27-218, Building 27, Seoul National University & +82-10-9982-8055 \\
1 Gwanak-ro, Gwanak-gu, Seoul 08826, South Korea & \email{kweon7182@snu.ac.kr}\\
\end{tabular}

%\section{Objective}

%Insert text here if you want to
%\begin{innerlist}
%\item More information and auxiliary documents can be found at\\\url{http://www.tedpavlic.com/facjobsearch/}
%\end{innerlist}

\section{Employment}
\href{https://math.snu.ac.kr}{\textbf{Seoul National University}}, Seoul, South Korea
\begin{outerlist}
\item[] BK Research Fellow \hfill Jun 2025--Present
\end{outerlist}

\vspace{.1in}
\href{https://www.math.uga.edu/}{\textbf{University of Georgia}}, Athens, GA, USA
\begin{outerlist}
\item[] Limited Term Assistant Professor \hfill Jan 2023--May 2025
\end{outerlist}

\vspace{.1in}
\href{http://www.mit.edu}{\textbf{Massachusetts Institute of Technology}}, Cambridge, MA, USA
\begin{outerlist}
\item[] Research Specialist Limited (One-year bridge postdoc) \hfill Aug 2021--Jul 2022
\end{outerlist}

\vspace{.1in}
\href{https://www.nims.re.kr/eng/}{\textbf{National Institute for Mathematical Sciences}}, Daejeon, South Korea
\begin{outerlist}
\item[] \textbf{Technical Research Personnel} \hfill Mar 2015--Apr 2018
\begin{innerlist}
\item Alternative military service
\item Managing a SageMath notebook server %\hfill Mar 2015--Dec 2015
\item Providing mathematical outreach programs  %\hfill Jan 2016--Apr 2018
%\item Encouraging people to engage in mathematical activities \\
%  particularly for underrepresented minorities
\end{innerlist}
\end{outerlist}

\section{Education}

\href{http://www.mit.edu}{\textbf{Massachusetts Institute of Technology}}, Cambridge, MA, USA
\begin{outerlist}
\item[] Ph.D. in Mathematics \hfill Sep 2018--Jun 2021 \\
(Prior enrollment before compulsory military service) \hfill Sep 2012--Aug 2014 \\
Advisor: Bjorn Poonen
\end{outerlist}

\vspace{.1in}
\href{http://www.postech.ac.kr}{\textbf{Pohang University of Science and Technology}}, Pohang, South Korea
\begin{outerlist}
\item[] B.S. in
        \href{http://math.postech.ac.kr/}
        {Mathematics} \hfill Mar 2008--Feb 2012\\
        B.S. in
        \href{https://cse.postech.ac.kr/}
        {Computer Science and Engineering}
        \begin{innerlist}
        \item {Double Major}
        \item {Summa Cum Laude} % (GPA: 4.15/4.30)}
        \end{innerlist}
\end{outerlist}

\section{Publications}
\href{https://arxiv.org/abs/2602.12718}{A Uniform Construction of Cohomology Theories of Varieties via Fundamental Groupoids}, \textit{preprint}\vspace{5pt}

\href{https://arxiv.org/abs/2601.16505}{Computing Picard Schemes} (with Madhavan Venkatesh),
\textit{submitted}\vspace{5pt}

\href{https://arxiv.org/abs/2511.00431}{Bounds on Torsion in Cohomology and an Effective Version of Deligne's Gcd Theorem} (with Madhavan Venkatesh),
\textit{to appear in Journal de Th\'eorie des Nombres de Bordeaux}\vspace{5pt}

\href{https://arxiv.org/abs/2301.02949}{Maximum Overlap Area of Several Convex Polygons Under Translation (with Honglin Zhu)}, \textit{The 36th Canadian Conference on Computational Geometry (CCCG 2024)}
\vspace{5pt}

\href{https://doi.org/10.4230/LIPIcs.SoCG.2023.61}{Maximum Overlap Area of a Convex Polyhedron and a Convex Polygon under Translation} (with Honglin Zhu),
\textit{The 39th International Symposium on Computational Geometry (SoCG 2023)}
\vspace{5pt}

\href{https://doi.org/10.1016/j.aim.2022.108687}{Bounds on the Torsion Subgroup Schemes of Néron-Severi Group Schemes, \textit{Advances in Mathematics 409 (2022): 108687.}}
\vspace{5pt}

\href{https://www.ams.org/journals/tran/0000-000-00/S0002-9947-2020-08203-8/S0002-9947-2020-08203-8.pdf}{Bounds on the Torsion Subgroups of Néron-Severi Groups, \textit{Trans. Amer. Math. Soc. 374 (2021), pp. 351--365.}}
\vspace{5pt}

\href{https://www.sciencedirect.com/science/article/pii/S0925772113000941}{Overlap of Convex Polytopes Under Rigid Motion} (with Hee-Kap Ahn, Siu-Wing Cheng, and Juyoung Yon), \textit{Computational Geometry} 47, no. 1 (2014): 15--24

%\halfblankline


\section{Awards and Honors}
AMS-Simons Travel Grants \hfill Jul 2023--Jun 2025

\vspace{5pt}

Scholarship for Doctoral Study, Samsung Foundation \hfill Sep 2012--Jun 2021

\vspace{5pt}

Scholarship, STX Foundation \hfill Mar 2011--Feb 2012

\vspace{5pt}

Governmental Scholarship, Korea Student Aid Foundation \hfill Mar 2008--Feb 2011

\vspace{5pt}

% Fellowship with Top Honor (5 times), POSTECH \hfill Mar 2008--Feb 2012

% Silver Award, National Collegiate Mathematics Competition, \hfill Nov 2010 \\
% the Korean Mathematical Society

% Five Golden Buttons, The Five-Gold-Button Math Competition, \hfill Sep 2008 \\
% POSTECH

6th Place, ACM-ICPC Korea Regional Programming Contest \hfill Nov 2010

\vspace{5pt}

4th Place, ACM-ICPC Korea Regional Programming Contest \hfill Nov 2009

\vspace{5pt}

5th Place, ACM-ICPC Korea Regional Programming Contest \hfill Nov 2008

\vspace{5pt}

% Top Prize, Programming Contest, POSTECH, Pohang, Korea \hfill May 2009


\section{Teaching}
\textbf{Calculus III for Engineering (MATH2500)} \hfill Jan 2025--May 2025 \\
Lecturer, University of Georgia\vspace{5pt}

\textbf{Complex Variables (MATH4150)} \hfill Aug 2024--Dec 2024 \\
Lecturer, University of Georgia\vspace{5pt}

\textbf{Precalculus (MATH1113)} \hfill Aug 2024--Dec 2024 \\
Lecturer, University of Georgia\vspace{5pt}

\textbf{Calculus II for Science and Engineering (MATH2260)} \hfill Jan 2024--Apr 2024 \\
Lecturer, University of Georgia\vspace{5pt}

\textbf{Calculus I for Science and Engineering (MATH2250)} \hfill Aug 2022--Dec 2023 \\
Lecturer, University of Georgia\vspace{5pt}

\textbf{Precalculus (MATH1113)} \hfill Jan 2022--May 2022 \\
Lecturer, University of Georgia\vspace{5pt}

\textbf{Multivariable Calculus (18.02)} \hfill Feb 2021--May 2021 \\
Recitation Leader, Massachusetts Institute of Technology\vspace{5pt}

\textbf{Calculus (18.01A/18.02A)} \hfill Sep 2020--Dec 2020 \\
Recitation Leader, Massachusetts Institute of Technology\vspace{5pt}

\textbf{Calculus of Several Variables (18.022)} \hfill Sep 2019--Dec 2019 \\
Recitation Leader, Massachusetts Institute of Technology\vspace{5pt}

\textbf{Differential Equations (MATH200)} \hfill Mar 2012--Jun 2012 \\
Recitation Leader, Pohang University of Science and Technology\vspace{5pt}

\textbf{Complex Variables (MATH310)} \hfill Aug 2011--Dec 2011\\
Recitation Leader, Pohang University of Science and Technology


\section{Research Talks}
International Conference for the 80th Anniversary of the \hfill Jun 2026\\
Korean Mathematical Society

\vspace{5pt}

Seminar, Soongsil University \hfill Mar 2026

\vspace{5pt}

KIAS HCMC Algebra Seminar, Korea Institute for Advanced Study \hfill Nov 2025

\vspace{5pt}

POSTECH-PMI Number Theory Seminar, POSTECH \hfill Nov 2025

\vspace{5pt}

Rookies Workshop 2025, Seoul National University \hfill Nov 2025

\vspace{5pt}

QSMS Summer Workshop 2025, Goseong \hfill Aug 2025

\vspace{5pt}

Rookies Pitch, Seoul National University \hfill Jun 2025

\vspace{5pt}

Number Theory Seminar, University of Georgia \hfill May 2025

\vspace{5pt}

Canadian Conference on Computational Geometry, Brock University \hfill Jul 2024

\vspace{5pt}

Palmetto Number Theory Series XXXVI, Clemson University \hfill Oct 2023

\vspace{5pt}

Number Theory Seminar, University of Georgia \hfill Apr 2023

\vspace{5pt}

Southern Regional Number Theory Conference, Louisiana State University \hfill Mar 2023

\vspace{5pt}

Number Theory Seminar, University of Georgia \hfill Aug 2022

\vspace{5pt}

Korea Institute for Advanced Study \hfill Oct 2022

\vspace{5pt}

POSTECH-PMI Number Theory Seminar, POSTECH \hfill Apr 2022

\vspace{5pt}

Johns Hopkins Junior Number Theory Days, Johns Hopkins University \hfill Dec 2021

\vspace{5pt}


2021 Korean Mathematical Society Annual Meeting \hfill Oct 2021

\vspace{5pt}

Ohio State Number Theory Seminar, Ohio State University \hfill Mar 2021

\vspace{5pt}

AMS Special Session, Spring Eastern Sectional Meeting, Brown University  \hfill Mar 2021

\vspace{5pt}

Binghamton Arithmetic Seminar, Binghamton University \hfill Mar 2021

\vspace{5pt}

Korea Institute for Advanced Study \hfill Jan 2020

\vspace{5pt}

% \sout{AMS Special Session, Spring Eastern Sectional Meeting, Tufts University} \hfill Mar 2020
% \\(Canceled due to COVID-19)
% \vspace{5pt}

\section{Expository Talks} % Expository
\textbf{Proof of the amplification principle 2} \hfill May 2022 \\
STAGE, Massachusetts Institute of Technology
\vspace{5pt}

\textbf{Alterations} \hfill Nov 2020 \\
STAGE, Massachusetts Institute of Technology
\vspace{5pt}

\textbf{Formal Groups as Functors} \hfill Sep 2019 \\
STAGE, Massachusetts Institute of Technology
\vspace{5pt}

\textbf{Division by 2 on Hyperelliptic Curves} \hfill Dec 2018 \\
Preprint Seminar, Massachusetts Institute of Technology
\vspace{5pt}

\textbf{Review of \'Etale Cohomology} \hfill Sep 2018 \\
STAGE, Massachusetts Institute of Technology
\vspace{5pt}

\textbf{On the Fary-Milnor Theorem, following the existing literature} \hfill Nov 2010 \\
SCV Seminar, Pohang University of Science and Technology

% \halfblankline

% \textbf{Co-researcher} \hfill Mar 2011--Dec 2011 \\
% Geometric Algorithms Lab, Department of Computer Science and Engineering, Pohang University of Science and Technology
% \begin{innerlist}
% \item Advisor: Prof. Hee-Kap Ahn
% \item Conducted research on approximation algorithms for finding a maximal overlap of two convex polyhedra under rigid motions.
% \end{innerlist}

\section{Academic Service}
Co-organizer of STAGE \hfill Spring 2022
\\(Seminar on Topics in Arithmetic, Geometry, Etc.), MIT \vspace{5pt}

Co-organizer of STAGE \hfill Fall 2021
\\(Seminar on Topics in Arithmetic, Geometry, Etc.), MIT \vspace{5pt}

Co-organizer of STAGE \hfill Fall 2020
\\(Seminar on Topics in Arithmetic, Geometry, Etc.), MIT

% \section{Activities}
% \textbf{Member of MARCUS} \hfill Mar 2008--Feb 2012
% \begin{innerlist}
% \item MARCUS is the undergraduate math club in POSTECH, and it has one-hour weekly math seminars.
% \end{innerlist}\newpage

% \textbf{Member of POSCAT} \hfill Apr 2008--Feb 2012
% \begin{innerlist}
% \item POSCAT is an academic club studying algorithms in POSTECH. We prepared for several programming contests, and developed artificial intelligence for annual AI competition between POSTECH and KAIST.
% \end{innerlist}

% \section{Technical \\ Skills}
% \textbf{Programming Languages}
% \begin{innerlist}
% \item C, C++, Python, OCaml
% \end{innerlist}

% \halfblankline

% \textbf{Mathematical Software}
% \begin{innerlist}
% \item Sage, Maple, Mathematica
% \end{innerlist}

% \halfblankline

% \textbf{Document Preparation System}
% \begin{innerlist}
% \item LaTeX
% \end{innerlist}

% \begin{comment}
% \end{comment}


\end{document}

%%%%%%%%%%%%%%%%%%%%%%%%%% End CV Document %%%%%%%%%%%%%%%%%%%%%%%%%%%%%

%----------------------------------------------------------------------%
% The following is copyright and licensing information for
% redistribution of this LaTeX source code; it also includes a liability
% statement. If this source code is not being redistributed to others,
% it may be omitted. It has no effect on the function of the above code.
%----------------------------------------------------------------------%
% Copyright (c) 2007, 2008, 2009, 2010, 2011 by Theodore P. Pavlic
%
% Unless otherwise expressly stated, this work is licensed under the
% Creative Commons Attribution-Noncommercial 3.0 United States License. To
% view a copy of this license, visit
% http://creativecommons.org/licenses/by-nc/3.0/us/ or send a letter to
% Creative Commons, 171 Second Street, Suite 300, San Francisco,
% California, 94105, USA.
%
% THE SOFTWARE IS PROVIDED "AS IS", WITHOUT WARRANTY OF ANY KIND, EXPRESS
% OR IMPLIED, INCLUDING BUT NOT LIMITED TO THE WARRANTIES OF
% MERCHANTABILITY, FITNESS FOR A PARTICULAR PURPOSE AND NONINFRINGEMENT.
% IN NO EVENT SHALL THE AUTHORS OR COPYRIGHT HOLDERS BE LIABLE FOR ANY
% CLAIM, DAMAGES OR OTHER LIABILITY, WHETHER IN AN ACTION OF CONTRACT,
% TORT OR OTHERWISE, ARISING FROM, OUT OF OR IN CONNECTION WITH THE
% SOFTWARE OR THE USE OR OTHER DEALINGS IN THE SOFTWARE.
%----------------------------------------------------------------------%
